\subsection{The fast and the hungry}

Isaiah asks the worshipper to look at the people who bear the cost of
his life. A fast could be kept in sackcloth while men still oppressed
their workers. God asks them instead to break the burdens they have laid
on others and to share their bread. “Deal thy bread to the hungry,” the
prophet says, before promising that light will rise like morning (Isa.
58:5--8). Suppose that a man has held back wages due on Friday. The person
at his door on Saturday asking for the money has already earned it.
Let him pay it first. He will then be free to consider what he can
give out of his own share. Isaiah puts the burden and the bread in
that order.

The fast also gives you something that can be handed over. On a day
each week when it is prudent, make a simpler meal and set aside what you
save. You could bring a bag of good food to the parish pantry or use that
money to help someone complete a week's shopping. If a change of diet is
imprudent, go without an optional purchase or give some of your free time
to someone who needs company or help. In his last words to Proba, Augustine
asks her household to bear with one another in charity while each prays,
fasts and keeps vigil according to what each can bear (Letter 130.31).

Jesus speaks of three things the Father sees done in secret. A hand
gives alms, a door closes for prayer, a man washes his face before
going out on a day when he is fasting (Matt. 6:1--6, 16--18). It is
easy for the person receiving a gift to become the subject of a story
we tell about ourselves. A parish can deliver groceries quietly.
A neighbour can take a dish to the door and come inside only if
invited. Christ entrusts both the giver's hidden sacrifice and the
recipient's dignity to the sight of his Father.

St John Chrysostom puts two cups before a congregation. One is the
gold cup someone hopes to give the church. The other is the cup of
cold water a thirsty man needs. The preacher takes his hearers back
to the Last Supper. The Lord who gave his disciples his body is the very one who
will ask at the judgment whether they fed him. Those who believe him
at the altar have every reason to believe him when he says he has
gone hungry. Chrysostom welcomes a gift for the altar from honest
earnings. He asks that the donor should feed the hungry Christ first
(Homily 50 on Matthew, \S\S4--5). As you work out what you can afford
to give this Christmas, find out who needs help even to reach the
table. In some families it is the price of the bus fare that keeps a
place empty.

A fast chosen for Advent is a voluntary practice. Under present
Latin Church law Fridays remain penitential days. Follow the
observance in force where you live. The laws that governed when
the 1962 Missal appeared and the laws in force now, including United
States practice, are set out in \emph{About this book}.

\subsection{Preparing for confession}

Find a time for confession before the last crowded days. Look up the hours
or arrange a time with a priest. If you are conscious of grave sin, seek
confession promptly. Sit quietly beforehand and ask with the psalmist for
a clean heart (Ps. 50:12). Look at what you have done since your last
confession. Ask where you have left prayer or a duty to God aside, where
you have spoken against your neighbour or withheld what you owed, and
where you have let a habit make a choice for you. Remember the good you
have received as well. It helps to know whom you are asking for pardon.

St Ignatius of Loyola begins his examination of a day by giving thanks
for benefits received. He then asks for light to recognize sin and goes
back over the time since rising, attending in turn to what he thought,
said and did. He ends by asking pardon, resolving with God's grace to
amend, and saying the Our Father (\emph{Spiritual Exercises}, Method
for Making the General Examen). If you can remember only the last angry
word of a quarrel, go back to breakfast. Someone did you a kindness
then. Later you mistook a remark for an insult. By evening you were
ready to give offence. Remembering the whole day may show you where
you can first ask for help tomorrow.

If you have been away from the sacrament a long time and have
forgotten how to begin, tell the priest just that. The Roman Catechism
calls sorrow for sin, with the purpose of leaving it behind, contrition.
That sorrow is something the will can do (Part II, Penance). You can
regret the harm you did and ask for the strength to change even on a
day when you feel very little. Consider where you can begin to act
differently when the same occasion comes.

In Jericho, Jesus stops under the tree where Zacchaeus has climbed to see
him. He tells the tax collector to come down because he means to stay in
his house. Zacchaeus makes haste and receives him with joy. In the house
he promises the poor half his goods and undertakes to pay back those he
has wronged. Jesus says salvation has come to that house (Luke 19:5--10).
The Lord asks to be received by a man whose dealings need to change.
He calls him out of the tree before the man has put his accounts in order.

When you confess, say plainly what needs forgiving. Ask for the grace to
turn from it. Through the priest's words of absolution Christ forgives,
exercising the mercy he entrusted to his apostles (John 20:22--23). Give thanks for the pardon you have received. Then think of the
person who still bears the cost of something you did. An overpayment can be returned. A story told at another
person's expense can be corrected with those who heard it. If you cannot
see what kind of repair would be just or possible, ask the confessor.
Do the penance he gives you. The Roman Catechism describes prayer,
fasting and alms as works of satisfaction that concern God, oneself
and one's neighbour (Part II, Penance). They are fitting ways for a
person who has received mercy to begin to live mercifully.
